%
%  Example solution set - shows how to include codes, results and figures
%
\documentclass[11pt]{article}
% 
\usepackage{amsmath}
\usepackage{amssymb}
\usepackage{graphicx}

\usepackage{calc}
\usepackage[margin=1.in]{geometry}

\usepackage[usenames]{color}

\usepackage{listings}
\lstset{
basicstyle=\footnotesize\ttfamily,
columns=flexible,
breaklines=true,
commentstyle=\color{red},
keywordstyle=\color{black}\bfseries,
keepspaces=true
}

\usepackage{fancyvrb}% fancy Verbatim

\newcommand{\Dp}{D_+}
\newcommand{\Dm}{D_-}
\newcommand{\Dz}{D_0}
\newcommand{\dx}{\Delta x}
\newcommand{\dy}{\Delta y}
\newcommand{\dt}{\Delta t}
\newcommand{\Nx}{N_x} 

\begin{document}

\begin{flushleft}
W.D. Henshaw \hfill \textbf{\large Math 4800 Sample Solution Set} \hfill ~~~~~~~~~~~~~~~~~
\end{flushleft}

\medskip
This file shows how to include code listings and output from Matlab into a LaTeX file.
It also shows how to include figures made in Matlab. 

% -------------------------
\bigskip
\noindent{\em Solution:}



% ----------------------------------------------------------------------------------------------
\bigskip
\noindent 1(a) The method-of-lines approximation is
\begin{align*}g
 &  \frac{d V_i(t)}{dt} = \kappa \Dp\Dm V_i(t) , \qquad i=1,2,3,\ldots,\Nx-1, \\
 &   V_0(t) = g_a(t), \qquad V_{\Nx}(t) = g_b(t), \\
 &   V_i(0) = \sin(k\pi x_i), \quad \qquad i=0,1,2,3,\ldots,\Nx
\end{align*}

\medskip
\noindent 1(b) Forward-Euler in time: 
\begin{align*}
 &  \frac{U_i^{n+1} - U_i^n}{\dt} = \kappa \Dp\Dm U_i^n , \qquad i=1,2,3,\ldots,\Nx-1, \\
 &   U_0^n = g_a(t^n), \qquad U_{\Nx}^n = g_b(t^n), \\
 &   U_i^0 = \sin(k\pi x_i), \quad \qquad i=0,1,2,3,\ldots,\Nx.
\end{align*}

\bigskip
\noindent 1(c) The exact solution is 
\begin{align*}
   u(x,t) = e^{ -  \kappa(k\pi)^2  \, t }\, \sin(k\pi x).
\end{align*}

\medskip
\noindent 1(d)  The matlab code is given below. Here are results for FE-CD2.

\medskip\noindent
Here is one way to include code output using the lstlisting environment: 
\begin{lstlisting}[frame=single,caption={Results from 1(d).},label=lst:FE-CD2]
>> sampleCode
 t=0.50 CPU=1.0e-02 Nx= 20 Nt=   44 dt=1.14e-02 max-Err=1.6e-03
 t=0.50 CPU=0.0e+00 Nx= 40 Nt=  178 dt=2.81e-03 max-Err=4.1e-04 ratio=4.01 rate=2.00
 t=0.50 CPU=3.0e-02 Nx= 80 Nt=  711 dt=7.03e-04 max-Err=1.0e-04 ratio=3.98 rate=1.99
 t=0.50 CPU=6.0e-02 Nx=160 Nt= 2844 dt=1.76e-04 max-Err=2.6e-05 ratio=4.00 rate=2.00
\end{lstlisting}

\medskip\noindent
Here is another way to include output using the Verbatim environment: 
\begin{Verbatim}[fontsize=\small]
>> sampleCode
 t=0.50 CPU=1.0e-02 Nx= 20 Nt=   44 dt=1.14e-02 max-Err=1.6e-03
 t=0.50 CPU=0.0e+00 Nx= 40 Nt=  178 dt=2.81e-03 max-Err=4.1e-04 ratio=4.01 rate=2.00
 t=0.50 CPU=3.0e-02 Nx= 80 Nt=  711 dt=7.03e-04 max-Err=1.0e-04 ratio=3.98 rate=1.99
 t=0.50 CPU=6.0e-02 Nx=160 Nt= 2844 dt=1.76e-04 max-Err=2.6e-05 ratio=4.00 rate=2.00
\end{Verbatim}

\lstinputlisting[language=Matlab, numbers=left, stepnumber=1, firstline=1,caption={heat.m, Matlab code for sample solution},label=code:heatEquation,frame=single]{sampleCode.m}

\begin{figure}[hbt]
\begin{center}
  \includegraphics[width=7.5cm]{heatEquation1d_FE} \quad
  \includegraphics[width=7.5cm]{heatEquation1d_FE_err}
  \caption{Results from problem 1. Left: computed and exact solutions at $t=1$. Right: errors at $t=1$.} \label{fig:heat}
\end{center}
\end{figure}


\end{document}
% =====================================================
% =====================================================
% =====================================================

